\chapter{Справочник Lua API}
\section{Модель выполнения}
app.lua возвращает таблицу M. M.loop() обязателен и вызывается многократно; M.init(), M.event(event\_type) и M.cleanup() необязательны. Без loop приложение отклоняется. При ошибках выполнения прошивка может перейти к резервному приложению.

Скрипт в sensor/*.lua возвращает таблицу sensor с sensor.get(index), возвращающим число. Прошивка делает sensor глобальным. apppath содержит путь к папке текущего приложения. Загружайте модули через require("module"); для необязательных модулей можно использовать pcall(require, "module").
\section{Ресурсы и отображение}
loadimg и loadfont возвращают дескриптор или nil; free освобождает загруженные изображения и шрифты. fill заполняет весь экран. update передаёт кадр, после чего очищает экран чёрным. drawpolygon требует чётное число координат, не менее четырёх элементов. Вызывайте update один раз на кадр и освобождайте ресурсы в cleanup.
\section{Другие модули}
system управляет задержкой, подсветкой, подтяжкой V2, выходом 5 В и сном, читает кнопки, АЦП и время. onewire ищет DS18B20 и читает температуру в градусах Цельсия. nvs использует пространство NVS app; get возвращает строку или nil, push — логическое значение. ble инициализирует BLE и предоставляет обороты, скорость и передачу.
\section{display}
\begin{itemize}
\item display.loadimg(path) -> img\_handle|nil
\item display.loadfont(path) -> font\_handle|nil
\item display.free(handle)
\item display.fill(color)
\item display.update()
\item display.drawimg(img, x, y)
\item display.string(font, x, y, text, color)
\item display.fillrect(x, y, width, height, color)
\item display.drawline(x1, y1, x2, y2, color)
\item display.drawcircle(x, y, r, color)
\item display.drawfillcircle(x, y, r, color)
\item display.drawroundrect(x1, y1, x2, y2, r, color)
\item display.drawarrow(x0, y0, x1, y1, w, color)
\item display.drawfillarrow(x0, y0, x1, y1, w, color)
\item display.drawpolygon(\{x1, y1, x2, y2, ...\}, color)
\item display.RED
\item display.GREEN
\item display.BLUE
\item display.WHITE
\item display.BLACK
\item display.YELLOW
\item display.CYAN
\item display.PURPLE
\item display.GRAY
\end{itemize}
\section{system}
\begin{itemize}
\item system.delay(ms)
\item system.backlight(level)
\item system.V2\_pullup(enable)
\item system.out5V(enable)
\item system.keys() -> int
\item system.adc\_ch(channel) -> number
\item system.get\_time() -> hour, minute, second
\item system.config\_sleep(enable)
\end{itemize}
\section{onewire}
\begin{itemize}
\item onewire.scan18b20() -> count
\item onewire.read18b20(index) -> temperature\_c
\end{itemize}
\section{nvs}
\begin{itemize}
\item nvs.get(key) -> string|nil
\item nvs.push(key, value\_string) -> boolean
\end{itemize}
\section{ble}
\begin{itemize}
\item ble.init() -> true
\item ble.get\_telemetry() -> rpm, speed, gear
\end{itemize}
\section{Примеры и рекомендации}
В примере вольтметра координаты находятся внутри экрана 160\,\texttimes\,80 пикселей, ширина полосы ограничена. Пример датчика показывает чтение АЦП и сглаживание, а не полный расчёт температуры NTC: замените иллюстративное преобразование реальной калибровкой. sensor.get должен работать быстро; храните калибровку через nvs и при необходимости преобразуйте строки в числа. Формируйте пути от apppath. Выполните поиск DS18B20 один раз, затем читайте датчики в цикле.
\begin{verbatim}
local M = {}
local font = display.loadfont(apppath .. "/../fonts/H24XB.FNT")

local function read_voltage()
return sensor.get(1)
end

function M.loop()
local v = read_voltage()
display.fill(display.BLACK)
display.string(font, 10, 30,
    string.format("%2.1f V", v), display.WHITE)
display.fillrect(10, 60,
    math.max(0, math.min(140, math.floor((v - 9.5) * 12))),
    8, display.GREEN)
display.update()
end

function M.cleanup()
display.free(font)
end

return M
\end{verbatim}
\begin{verbatim}
sensor = {}

local filt = 0.0

function sensor.get(i)
i = i or 0
system.V2_pullup(1)

local raw = system.adc_ch(i)
filt = filt + (raw - filt) / 10.0

local value = filt * 1.0
return value
end

return sensor
\end{verbatim}
